\chapter{Symmetric Spaces and Holonomy}
\label{chap:pet-symmetric-spaces-and-holonomy}

Symmetric spaces are characterized locally by parallel curvature and admit a
Lie-algebraic description at each point. Holonomy then controls metric
splittings, parallel tensors, and the irreducible curvature alternatives.

\section{Symmetric Spaces}

\subsection{The Homogeneous Description}

Point symmetries imply transitivity and supply the basic compact and
noncompact families.

\begin{definition}[symmetric space]
  \label{def:pet-ch10-symmetric-space}
  \lean{PetersenLib.IsSymmetricSpace}
  \source{petersen:page-0379}
  A Riemannian manifold $(M,g)$ is \emph{symmetric} if every $p\in M$ admits
  $A_p\in\mathrm{Iso}_p$ with $D_pA_p=-I$. This \emph{point symmetry} reverses
  each geodesic through $p$:
  $A_p(c(t))=c(-t)$ whenever $c(0)=p$.
\end{definition}

\begin{proposition}[symmetric spaces are homogeneous, hence complete]
  \label{prop:pet-ch10-symmetric-homogeneous}
  \lean{PetersenLib.symmetricSpace_isHomogeneous}
  \source{petersen:page-0379}
  \uses{def:pet-ch10-symmetric-space}
  A symmetric space is homogeneous, and consequently complete.
\end{proposition}
\begin{proof}
  \uses{prop:pet-ch10-symmetric-homogeneous}
  \sketch
  The symmetry at the midpoint of a geodesic segment interchanges its
  endpoints. Composing such symmetries along a broken geodesic sends any point
  to any other, so the isometry group is transitive. The maximal existence
  time of a geodesic is therefore independent of its initial point; local
  existence then extends every geodesic for all time.
\end{proof}

\begin{remark}[a symmetry at one point suffices; biinvariant Lie groups]
  \label{rem:pet-ch10-single-point-symmetry}
  \lean{PetersenLib.symmetricSpace_of_symmetry_at_one_point,PetersenLib.biinvariantLieGroupSymmetricSpace}
  \source{petersen:page-0379}
  \uses{def:pet-ch10-symmetric-space}
  On a homogeneous space, one point symmetry determines all others: if
  $q=h\cdot p$, then $A_q=hA_ph^{-1}$. Consequently every Lie group with a
  biinvariant metric is symmetric, because inversion is the point symmetry at
  the identity.
\end{remark}

\begin{definition}[classical symmetric-space families]
  \label{def:pet-ch10-symmetric-space-tables}
  \lean{PetersenLib.symmetricSpaceTables}
  \source{petersen:page-0379,petersen:page-0380}
  \uses{rem:pet-ch10-single-point-symmetry}
  The following compact symmetric spaces and their noncompact duals have the
  indicated ranks and dimensions.
  \begin{center}
    \begin{tabular}{c|c|c}
      \multicolumn{3}{c}{\textbf{Table 10.1 Compact Groups}}\\
      group & rank & dim\\\hline
      $\mathrm{SU}(n+1)$ & $n$ & $n(n+2)$\\
      $\mathrm{SO}(2n+1)$ & $n$ & $n(2n+1)$\\
      $\mathrm{Sp}(n)$ & $n$ & $n(2n+1)$\\
      $\mathrm{SO}(2n)$ & $n$ & $n(2n-1)$
    \end{tabular}
    \qquad
    \begin{tabular}{c|c|c}
      \multicolumn{3}{c}{\textbf{Table 10.2 Noncompact Analogues}}\\
      (complexified group)/group & rank & dim\\\hline
      $\mathrm{SL}(n{+}1,\mathbb C)/\mathrm{SU}(n{+}1)$ & $n$ & $n(n+2)$\\
      $\mathrm{SO}(2n{+}1,\mathbb C)/\mathrm{SO}(2n{+}1)$ & $n$ & $n(2n+1)$\\
      $\mathrm{Sp}(n,\mathbb C)/\mathrm{Sp}(n)$ & $n$ & $n(2n+1)$\\
      $\mathrm{SO}(2n,\mathbb C)/\mathrm{SO}(2n)$ & $n$ & $n(2n-1)$
    \end{tabular}
  \end{center}
  \begin{center}
    \begin{tabular}{l|l|c|c|l}
      \multicolumn{5}{c}{\textbf{Table 10.3 Compact Homogeneous Spaces}}\\
      $\mathrm{Iso}$ & $\mathrm{Iso}_p$ & dim & rank & description\\\hline
      $\mathrm{SO}(n{+}1)$ & $\mathrm{SO}(n)$ & $n$ & $1$ & sphere\\
      $\mathrm O(n{+}1)$ & $\mathrm O(n)\times\{1,-1\}$ & $n$ & $1$ & $\mathbb{RP}^n$\\
      $\mathrm U(n{+}1)$ & $\mathrm U(n)\times\mathrm U(1)$ & $2n$ & $1$ & $\mathbb{CP}^n$\\
      $\mathrm{Sp}(n{+}1)$ & $\mathrm{Sp}(n)\times\mathrm{Sp}(1)$ & $4n$ & $1$ & $\mathbb{HP}^n$\\
      $\mathrm F_4$ & $\mathrm{Spin}(9)$ & $16$ & $1$ & $\mathbb{OP}^2$\\
      $\mathrm{SO}(p{+}q)$ & $\mathrm{SO}(p)\times\mathrm{SO}(q)$ & $pq$ & $\min(p,q)$ & real Grassmannian\\
      $\mathrm{SU}(p{+}q)$ & $\mathrm S(\mathrm U(p)\times\mathrm U(q))$ & $2pq$ & $\min(p,q)$ & complex Grassmannian\\
      $\mathrm{SU}(n)$ & $\mathrm{SO}(n)$ & $\tfrac{n^2-1}2$ & $n-1$ & $\mathbb R^n$'s in $\mathbb C^n$
    \end{tabular}
  \end{center}
  \begin{center}
    \begin{tabular}{l|l|c|c|l}
      \multicolumn{5}{c}{\textbf{Table 10.4 Noncompact Homogeneous Spaces}}\\
      $\mathrm{Iso}$ & $\mathrm{Iso}_p$ & dim & rank & description\\\hline
      $\mathrm{SO}(n,1)$ & $\mathrm{SO}(n)$ & $n$ & $1$ & hyperbolic space\\
      $\mathrm O(n,1)$ & $\mathrm O(n)\times\{1,-1\}$ & $n$ & $1$ & hyperbolic $\mathbb{RP}^n$\\
      $\mathrm U(n,1)$ & $\mathrm U(n)\times\mathrm U(1)$ & $2n$ & $1$ & hyperbolic $\mathbb{CP}^n$\\
      $\mathrm{Sp}(n,1)$ & $\mathrm{Sp}(n)\times\mathrm{Sp}(1)$ & $4n$ & $1$ & hyperbolic $\mathbb{HP}^n$\\
      $\mathrm F_4^{-20}$ & $\mathrm{Spin}(9)$ & $16$ & $1$ & hyperbolic $\mathbb{OP}^2$\\
      $\mathrm{SO}(p,q)$ & $\mathrm{SO}(p)\times\mathrm{SO}(q)$ & $pq$ & $\min(p,q)$ & hyperbolic Grassmannian\\
      $\mathrm{SU}(p,q)$ & $\mathrm S(\mathrm U(p)\times\mathrm U(q))$ & $2pq$ & $\min(p,q)$ & complex hyperbolic Grassmannian\\
      $\mathrm{SL}(n,\mathbb R)$ & $\mathrm{SO}(n)$ & $\tfrac{n^2-1}2$ & $n-1$ & Euclidean structures on $\mathbb R^n$
    \end{tabular}
  \end{center}
  Here $\mathrm{Spin}(n)$ is the universal double cover of $\mathrm{SO}(n)$ for
  $n>2$, and
  $\mathrm{SO}(2)=\mathrm U(1)$, $\mathrm{Spin}(3)=\mathrm{SU}(2)=\mathrm{Sp}(1)$,
  $\mathrm{Spin}(4)=\mathrm{Spin}(3)\times\mathrm{Spin}(3)$. O'Neill's formula
  gives $\sec\ge0$ for the compact examples. Each compact or noncompact
  projective space is quarter-pinched: the ratio of extremal sectional
  curvatures is $\tfrac14$.
\end{definition}

\begin{definition}[rank of a geodesic and of a Riemannian manifold; CROSSes]
  \label{def:pet-ch10-rank}
  \lean{PetersenLib.geodesicRank,PetersenLib.manifoldRank}
  \source{petersen:page-0381}
  For a geodesic $c:\mathbb R\to M$, its \emph{rank} is the dimension of
  \[
    \{E:E\text{ is parallel along }c,
      \ R(E(t),\dot c(t))\dot c(t)=0\text{ for every }t\}.
  \]
  This is at least one because $E=\dot c$ is admissible. The rank of $(M,g)$
  is the minimum over all geodesics. Products have rank at least two. The
  rank-one spaces in \cref{def:pet-ch10-symmetric-space-tables} are the
  compact and noncompact rank-one symmetric spaces; the compact ones are
  called \emph{CROSSes}.
\end{definition}

\subsection{Isometries and Parallel Curvature}

The invariance of $\nabla R$ under a point symmetry gives, at its fixed point,
\[
  -(\nabla_XR)(Y,Z,W)
  =(\nabla_{-X}R)(-Y,-Z,-W)=(\nabla_XR)(Y,Z,W).
\]
Thus $\nabla R=0$ there and, by
\cref{prop:pet-ch10-symmetric-homogeneous}, throughout a symmetric space.

\begin{definition}[locally symmetric space]
  \label{def:pet-ch10-locally-symmetric-space}
  \lean{PetersenLib.IsLocallySymmetric}
  \source{petersen:page-0382}
  A Riemannian manifold is \emph{locally symmetric} when its curvature tensor
  is parallel, $\nabla R=0$.
\end{definition}

\begin{theorem}[local symmetries from parallel curvature]
  \label{thm:pet-ch10-cartan-local-symmetry}
  \lean{PetersenLib.cartanLocalSymmetryExistence}
  \source{petersen:page-0381,petersen:page-0382}
  \uses{def:pet-ch10-locally-symmetric-space, def:pet-ch10-symmetric-space}
  \dcref{ch10:10.1.1}
  Every $p$ in a locally symmetric manifold has a neighborhood on which an
  isometry $A_p$ satisfies $D_pA_p=-I$. If the manifold is also simply
  connected and complete, $A_p$ extends globally and the manifold is
  symmetric.
\end{theorem}
\begin{proof}
  \uses{thm:pet-ch10-cartan-local-symmetry,thm:pet-ch10-cartan-isometry-extension}
  \sketch
  On a normal ball define $A_p=\exp_p\circ(-I)\circ\exp_p^{-1}$. Radial
  curvatures determine the metric in exponential coordinates. For opposite
  radial geodesics their initial curvature endomorphisms agree, and
  $\nabla_{\partial_r}R=0$ keeps them equal; hence $A_p$ is a local isometry.
  The continuation argument used in
  \cref{thm:pet-ch10-cartan-isometry-extension} extends it along every curve.
  Completeness prevents finite-time termination, and simple connectivity makes
  the extension path-independent. Its differential at $p$ remains $-I$.
\end{proof}

\begin{remark}[left-invariant metrics need not be locally symmetric]
  \label{rem:pet-ch10-non-locally-symmetric-examples}
  \lean{PetersenLib.bergerSphereNotLocallySymmetric,PetersenLib.heisenbergNotLocallySymmetric}
  \source{petersen:page-0382}
  \uses{def:pet-ch10-locally-symmetric-space}
  For $\varepsilon\ne1$, Berger spheres and the Heisenberg group with their
  standard left-invariant metrics have neither parallel curvature nor, in
  dimension three, parallel Ricci tensor.
\end{remark}

\begin{theorem}[extending a curvature-preserving linear isometry]
  \label{thm:pet-ch10-cartan-isometry-extension}
  \lean{PetersenLib.cartanIsometryExtension}
  \source{petersen:page-0382}
  \uses{def:pet-ch10-symmetric-space, def:pet-ch10-locally-symmetric-space}
  \dcref{ch10:10.1.2}
  Let $(M,g)$ be simply connected and symmetric, and let $(N,\bar g)$ be a
  complete locally symmetric space of the same dimension. For $p\in M$ and
  $q\in N$, suppose the linear isometry $L:T_pM\to T_qN$ satisfies
  \[
    L(R^\sharp(x,y)z)=\bar R^\sharp(Lx,Ly)Lz
    \quad\text{for all }x,y,z\in T_pM.
  \]
  Then there is a unique Riemannian isometry $F:M\to N$ with $D_pF=L$.
\end{theorem}
\begin{proof}
  \uses{thm:pet-ch10-cartan-isometry-extension}
  \sketch
  On normal neighborhoods the only possible extension is
  $F=\exp_q\circ L\circ\exp_p^{-1}$. Parallel translation of $L$ identifies the
  curvature tensors along corresponding radial geodesics because both are
  parallel and agree initially. The radial metric equations then show that
  $F$ is a local isometry, as in
  \cref{thm:pet-ch10-cartan-local-symmetry}. Analytic continuation is
  single-valued by simple connectivity of $M$ and remains defined by
  completeness of $N$. The first jet at one point also gives uniqueness; the
  parallel-curvature input is \cref{def:pet-ch10-locally-symmetric-space}.
\end{proof}

\subsection{The Lie Algebra Description}

Fix a symmetric space $(M,g)$, a point $p$, and its point symmetry $A_p$. Write
$\mathrm{iso}=\mathrm{iso}(M,g)$,
$\mathrm{iso}_p=\{X\in\mathrm{iso}:X_p=0\}$, and define
\[
  \sigma_p=(A_p)_*:\mathrm{iso}\longrightarrow\mathrm{iso},\qquad
  \sigma_p(X)=DA_p\circ X\circ A_p^{-1}.
\]

\begin{proposition}[the involution $\sigma_p$ on $\mathrm{iso}$]
  \label{prop:pet-ch10-sigma-p-involution}
  \lean{PetersenLib.sigmaPInvolution}
  \source{petersen:page-0383,petersen:page-0384}
  \dcref{ch10:10.1.3}
  The map $\sigma_p=(A_p)_*$ is an involution of $\mathrm{iso}$: $\sigma_p^2=\mathrm{id}$.
  Its $1$-eigenspace is $\mathrm{iso}_p$, and its $(-1)$-eigenspace is
  $\{X\in\mathrm{iso}\mid(\nabla X)|_p=0\}$.
\end{proposition}
\begin{proof}
  \uses{prop:pet-ch10-sigma-p-involution}
  \sketch
  From $A_p^2=I$ one obtains $\sigma_p^2=I$. Naturality of the connection and
  $D_pA_p=-I$ give, for every Killing field $X$,
  \[
    \sigma_p(X)_p=-X_p,\qquad
    (\nabla\sigma_p(X))_p=(\nabla X)_p.
  \]
  Thus a $1$-eigenvector vanishes at $p$, while a $(-1)$-eigenvector has zero
  covariant derivative there. Conversely, the displayed identities make
  $X$ and $\sigma_pX$, or $-X$ and $\sigma_pX$, have the same first jet.
  Uniqueness of a Killing field from its first jet proves both reverse
  inclusions.
\end{proof}

The first-jet exact sequence and homogeneity identify the evaluation image of
$\mathrm{iso}$ with $T_pM$. By \cref{prop:pet-ch10-sigma-p-involution}, evaluation
restricts to an isomorphism from the $(-1)$-eigenspace of $\sigma_p$ onto
$T_pM$.

\begin{definition}[the subspaces $t_p$ and $s_p$]
  \label{def:pet-ch10-tp-decomposition}
  \lean{PetersenLib.tpSubspace,PetersenLib.spSubspace}
  \source{petersen:page-0384}
  \uses{prop:pet-ch10-sigma-p-involution}
  Let $t_p$ be the $(-1)$-eigenspace of $\sigma_p$. Then
  \cref{prop:pet-ch10-sigma-p-involution} gives
  \[
    t_p=\{X\in\mathrm{iso}:(\nabla X)_p=0\},\qquad
    t_p\xrightarrow{\ X\mapsto X_p\ }T_pM
  \]
  as a linear isomorphism. Hence $\mathrm{iso}=t_p\oplus\mathrm{iso}_p$, and
  \[
    \mathrm{iso}\simeq T_pM\oplus s_p\subset T_pM\oplus\mathfrak{so}(T_pM),
    \qquad
    s_p=\{(\nabla X)|_p\in\mathfrak{so}(T_pM)\mid X\in\mathrm{iso}_p\}.
  \]
\end{definition}

\begin{proposition}[the bracket of $\mathrm{iso}$ from the curvature]
  \label{prop:pet-ch10-lie-bracket-curvature}
  \lean{PetersenLib.symmetricSpaceLieBracketCurvature}
  \source{petersen:page-0384,petersen:page-0385}
  \uses{def:pet-ch10-tp-decomposition}
  \dcref{ch10:10.1.4}
  Under \cref{def:pet-ch10-tp-decomposition}, the bracket is determined by
  $R_p$:
  \begin{enumerate}
    \item If $X,Y\in T_pM$, then $[X,Y]=R(X,Y)\in s_p$.
    \item If $X\in T_pM$ and $S\in s_p$, then $[X,S]=-S(X)\in T_pM$.
    \item If $S,T\in s_p$, then $[S,T]=-(S\circ T-T\circ S)\in s_p$.
  \end{enumerate}
\end{proposition}
\begin{proof}
  \uses{prop:pet-ch10-lie-bracket-curvature}
  \sketch
  Use torsion-freeness together with
  \[
    [\nabla X,\nabla Y](V)+\nabla_V[X,Y]=R(X,Y)V.
  \]
  If $X,Y\in t_p$, their first derivatives vanish at $p$, so $[X,Y]_p=0$
  and $(\nabla[X,Y])_p=R(X_p,Y_p)$, proving the first formula. If
  $X\in t_p$ and $S=(\nabla Y)_p$ for $Y\in\mathrm{iso}_p$, evaluation of
  $[X,Y]=\nabla_YX-\nabla_XY$ gives $[X,S]=-S(X)$. Finally, for isotropy
  fields with derivatives $S,T$, both values vanish at $p$; the displayed
  curvature identity reduces to $(\nabla[X,Y])_p=-(S\circ T-T\circ S)$.
\end{proof}

\begin{definition}[adjoint action and Killing form]
  \label{def:pet-ch10-killing-form}
  \lean{PetersenLib.adjointAction,PetersenLib.killingForm}
  \source{petersen:page-0385}
  For a finite-dimensional Lie algebra $\mathfrak g$, define
  \[
    \mathrm{ad}_X(Y)=[X,Y],\qquad
    B(X,Y)=\operatorname{tr}(\mathrm{ad}_X\mathrm{ad}_Y).
  \]
  The Killing form $B$ is symmetric and invariant:
  $B(\mathrm{ad}_XY,Z)+B(Y,\mathrm{ad}_XZ)=0$.
\end{definition}

\begin{remark}[$\mathrm{ad}_S$ is skew and $B\le0$ on $s_p$]
  \label{rem:pet-ch10-adjoint-skew-symmetric}
  \lean{PetersenLib.adjointOnSpNegativeSemidefinite}
  \source{petersen:page-0385}
  \uses{prop:pet-ch10-lie-bracket-curvature, def:pet-ch10-killing-form}
  \dcref{ch10:10.1.5}
  Give $T_pM\oplus s_p$ the orthogonal-sum metric, using
  $\langle S_1,S_2\rangle=-\operatorname{tr}(S_1S_2)$ on $s_p$. For $S\in s_p$,
  $\mathrm{ad}_S$ is skew-adjoint. Consequently
  \[
    B(S,S)=\operatorname{tr}(\mathrm{ad}_S\circ\mathrm{ad}_S)=-\operatorname{tr}\big(\mathrm{ad}_S\circ(\mathrm{ad}_S)^*\big)\le0,
  \]
  with equality only if $S=0$.
\end{remark}
\begin{proof}
  \uses{rem:pet-ch10-adjoint-skew-symmetric}
  \sketch
  Skew-adjointness yields
  $B(S,S)=-\operatorname{tr}(\mathrm{ad}_S\mathrm{ad}_S^*)=-\|\mathrm{ad}_S\|^2$.
  Equality therefore implies $\mathrm{ad}_S=0$. Its action on $T_pM$ is $S$
  by \cref{prop:pet-ch10-lie-bracket-curvature}, so equality forces $S=0$.
\end{proof}

\begin{theorem}[curvature from the bracket]
  \label{thm:pet-ch10-curvature-bracket-formula}
  \lean{PetersenLib.curvatureFromBracket,PetersenLib.ricciFromKillingForm}
  \source{petersen:page-0385,petersen:page-0386,petersen:page-0387}
  \uses{prop:pet-ch10-lie-bracket-curvature, def:pet-ch10-killing-form}
  \dcref{ch10:10.1.6}
  If $X,Y,Z\in t_p$, then at $p$, $R(X,Y)Z=[Z,[X,Y]]$; equivalently, under
  $\mathrm{iso}\simeq T_pM\oplus s_p$,
  \[
    R(X,Y)Z=-\mathrm{ad}_Z\circ\mathrm{ad}_Y(X),
    \qquad
    \operatorname{Ric}(Y,Z)=-\tfrac12B(Z,Y).
  \]
  Moreover, if $\operatorname{Ric}=\lambda g$ for a scalar $\lambda\ne0$, the curvature
  operator has the same sign as $\lambda$.
\end{theorem}
\begin{proof}
  \uses{thm:pet-ch10-curvature-bracket-formula,rem:pet-ch10-adjoint-skew-symmetric}
  \sketch
  For Killing fields, $\nabla^2_{A,B}C=-R(C,A)B$. Apply this twice to
  $X,Y,Z\in t_p$ and use Bianchi to obtain
  \[
    \nabla_Z[X,Y]=R(X,Y)Z=[Z,[X,Y]].
  \]
  The bracket rules in \cref{prop:pet-ch10-lie-bracket-curvature} give the
  adjoint formula. Since $\mathrm{ad}_Y$ and $\mathrm{ad}_Z$ interchange the two
  summands of $T_pM\oplus s_p$, cyclicity of trace gives
  \[
    B(Z,Y)=2\operatorname{tr}
      (\mathrm{ad}_Z\mathrm{ad}_Y|_{T_pM}),
    \qquad \operatorname{Ric}(Z,Y)=-\tfrac12B(Z,Y).
  \]
  If $\operatorname{Ric}=\lambda g$, invariance of $B$ converts the curvature
  quadratic form on $\sum_iX_i\wedge Y_i$ into
  \[
    -\frac1{2\lambda}
    B\left(\sum_i\mathrm{ad}_{X_i}(Y_i),
            \sum_i\mathrm{ad}_{X_i}(Y_i)\right).
  \]
  The vector inside $B$ belongs to $s_p$, where
  \cref{rem:pet-ch10-adjoint-skew-symmetric} makes $B$ nonpositive. The
  curvature operator consequently has the sign of $\lambda$.
\end{proof}

\begin{definition}[the abstract bracket on $T_pM\oplus\mathfrak{so}(T_pM)$]
  \label{def:pet-ch10-abstract-symmetric-bracket}
  \lean{PetersenLib.abstractSymmetricSpaceBracket}
  \source{petersen:page-0387}
  \uses{thm:pet-ch10-curvature-bracket-formula}
  For any Riemannian manifold and $p\in M$, put
  $\mathfrak N_p=T_pM\oplus\mathfrak{so}(T_pM)$ and define
  \[
    [X,Y]=R(X,Y)\in\mathfrak{so}(T_pM)\quad(X,Y\in T_pM),
  \]
  \[
    -[S,X]=[X,S]=S(X)\in T_pM\quad(X\in T_pM,\ S\in\mathfrak{so}(T_pM)),
  \]
  \[
    [S,S']=-(S\circ S'-S'\circ S)\in\mathfrak{so}(T_pM)\quad(S,S'\in\mathfrak{so}(T_pM)).
  \]
  Jacobi is automatic for triples containing at most one tangent vector. For
  three tangent vectors it is the first Bianchi identity. For $X,Y\in T_pM$
  and $S\in\mathfrak{so}(T_pM)$, the remaining Jacobi expression is
  $-S\circ R_{X,Y}+R_{X,Y}\circ S+R_{S(X),Y}+R_{X,S(Y)}$ for $X,Y\in T_pM$,
  and need not vanish. Explicitly, the all-tangent relation is
  \[
    R(X,Y)Z+R(Z,X)Y+R(Y,Z)X
    =[Z,[X,Y]]+[Y,[Z,X]]+[X,[Y,Z]]=0.
  \]
\end{definition}

\begin{corollary}[membership criterion for $s_p$]
  \label{cor:pet-ch10-so-membership-criterion}
  \lean{PetersenLib.spMembershipCriterion}
  \source{petersen:page-0388}
  \uses{def:pet-ch10-tp-decomposition}
  \dcref{ch10:10.1.7}
  Let $(M,g)$ be a simply connected symmetric space. Then $S\in\mathfrak{so}(T_pM)$
  lies in $s_p$ if and only if, for all $X,Y\in T_pM$,
  \[
    -S\circ R_{X,Y}+R_{X,Y}\circ S+R_{S(X),Y}+R_{X,S(Y)}=0.
  \]
\end{corollary}
\begin{proof}
  \uses{cor:pet-ch10-so-membership-criterion,thm:pet-ch10-cartan-isometry-extension}
  \sketch
  For $Z\in\mathrm{iso}_p$ and $S=(\nabla Z)_p$, the identity
  $\mathcal L_ZR=0$ evaluated at $p$ is
  \[
    0=-S(R_{X,Y}V)+R_{X,Y}S(V)+R_{S(X),Y}V+R_{X,S(Y)}V,
  \]
  which proves necessity. Conversely, the asserted relation says that the
  one-parameter group generated by $S$ preserves $R_p$. By
  \cref{thm:pet-ch10-cartan-isometry-extension}, these linear maps extend to
  global isometries fixing $p$. Their infinitesimal generator lies in
  $\mathrm{iso}_p$ and has derivative $S$.
\end{proof}

\begin{corollary}[$T_pM\oplus s_p$ is the maximal Lie subalgebra]
  \label{cor:pet-ch10-maximal-subalgebra}
  \lean{PetersenLib.tpSpMaximalSubalgebra}
  \source{petersen:page-0388}
  \uses{def:pet-ch10-abstract-symmetric-bracket, cor:pet-ch10-so-membership-criterion}
  \dcref{ch10:10.1.8}
  Let $\mathfrak r_p\subset s_p$ be the Lie algebra generated by the curvature
  endomorphisms $R_{X,Y}$, $X,Y\in T_pM$. For a simply connected symmetric
  space, the bracket in
  \cref{def:pet-ch10-abstract-symmetric-bracket} makes $T_pM\oplus s_p$ the
  maximal Lie algebra of this form, with subalgebra
  $c_p=T_pM\oplus\mathfrak r_p$. Moreover,
  \[
    c_p\subset T_pM\oplus s_p
      \subset\mathfrak N_{\mathfrak N_p}(c_p).
  \]
  The canonical involution on $T_pM\oplus s_p$ has $T_pM$ and $s_p$ as its
  $(-1)$- and $1$-eigenspaces; its restriction to $c_p$ has eigenspaces
  $T_pM$ and $\mathfrak r_p$.
\end{corollary}
\begin{proof}
  \uses{cor:pet-ch10-maximal-subalgebra}
  \sketch
  The only nonautomatic Jacobi relation for the bracket in
  \cref{def:pet-ch10-abstract-symmetric-bracket} is the mixed relation. By
  \cref{cor:pet-ch10-so-membership-criterion}, every
  $\mathfrak t\subset\mathfrak{so}(T_pM)$ for which
  $T_pM\oplus\mathfrak t$ is a Lie algebra satisfies $\mathfrak t\subset s_p$.
  This proves maximality. Since $[X,X']=R_{X,X'}$, the space
  $T_pM\oplus\mathfrak r_p$ is a subalgebra. For $S\in s_p$, the identity in
  \cref{cor:pet-ch10-so-membership-criterion} shows that
  $[S,R_{X,Y}]$ is a linear combination of curvature endomorphisms, while
  $[S,X]=-S(X)\in T_pM$. Thus $T_pM\oplus s_p$ normalizes $c_p$. The two
  summands have signs $-1$ and $1$ under the canonical involution, giving the
  stated eigenspaces and their restrictions to $c_p$.
\end{proof}

\begin{remark}[reconstructing symmetric spaces from an involutive Lie algebra]
  \label{rem:pet-ch10-symmetric-space-from-involution}
  \lean{PetersenLib.symmetricSpaceFromInvolutiveLieAlgebra}
  \source{petersen:page-0388,petersen:page-0389}
  \uses{cor:pet-ch10-maximal-subalgebra}
  Let $\sigma$ be an involution of a Lie algebra $\mathfrak g$, with
  $\mathfrak g=\mathfrak t\oplus\mathfrak k$ its $(-1)$- and $1$-eigenspaces.
  Then
  \[
    [\mathfrak k,\mathfrak k]\subset\mathfrak k,\qquad
    [\mathfrak k,\mathfrak t]\subset\mathfrak t,\qquad
    [\mathfrak t,\mathfrak t]\subset\mathfrak k.
  \]
  Suppose a connected compact group $K$ with Lie algebra $\mathfrak k$
  integrates the bracket action on $\mathfrak t$; this is automatic when $K$
  is simply connected. Averaging gives a $K$-invariant Euclidean metric on
  $\mathfrak t$. If $K\subset G$ and $\operatorname{Lie}(G)=\mathfrak g$, this
  metric defines a Riemannian quotient $G/K$. For simply connected $G$,
  $\sigma$ integrates to an involution of $G$ fixing $K$, and its action at
  $eK$ is the point symmetry.

  The relevant homotopy sequence is
  \[
    \pi_1(K)\to\pi_1(G)\to\pi_1(G/K)\to\pi_0(K)\to\pi_0(G)\to\pi_0(G/K)\to1,
  \]
  so $G/K$ is connected and simply connected if
  $\pi_0(K)\to\pi_0(G)$ is an isomorphism and
  $\pi_1(K)\to\pi_1(G)$ is surjective. The involution is essential:
  $-\mathrm{id}$ is an anti-automorphism, not an automorphism, on a nonabelian
  Lie algebra.
\end{remark}

\section{Examples of Symmetric Spaces}

For $X,Y\in\operatorname{Mat}_{k\times l}(\mathbb F)$, where
$\mathbb F\in\{\mathbb R,\mathbb C\}$, set
\[
  \langle X,Y\rangle=\operatorname{tr}(X^\top Y)
  =\overline{\operatorname{tr}(XY^*)}.
\]
This inner product is invariant under $X\mapsto AXB^\top$ for $A\in O(k)$ and
$B\in O(l)$.

\begin{example}[the compact Grassmannian]
  \label{ex:pet-ch10-compact-grassmannian}
  \lean{PetersenLib.compactGrassmannianSymmetricSpace}
  \source{petersen:page-0390,petersen:page-0391,petersen:page-0392}
  \uses{def:pet-ch10-tp-decomposition, def:pet-ch10-killing-form, thm:pet-ch10-curvature-bracket-formula, cor:pet-ch10-maximal-subalgebra}
  \dcref{ch10:10.2.1}
  Let $k,l\ge1$, $k+l\ge3$, and let
  $M=\widetilde G_k(\mathbb R^{k+l})$ be the oriented Grassmannian based at
  $p=\mathbb R^k$. Then
  \[
    M=O(k+l)/(\mathrm{SO}(k)\times O(l)),\qquad
    T_pM\simeq\operatorname{Mat}_{k\times l}(\mathbb R),
  \]
  with isotropy action $(A,B)\cdot X=AXB^\top$. The element $(I_k,-I_l)$
  induces the point symmetry, and
  \[
    \mathfrak{so}(k+l)=t_p\oplus\mathfrak{so}(k)\oplus\mathfrak{so}(l),
    \qquad
    t_p=\Big\{\begin{pmatrix}0&B\\-B^\top&0\end{pmatrix}\mid B\in\operatorname{Mat}_{k\times l}\Big\},
  \]
  On $t_p$,
  $\langle X,Y\rangle=\operatorname{tr}(X^\top Y)=-\operatorname{tr}(XY)
  =-(k+l-2)^{-1}B(X,Y)$, whence
  \[
    \operatorname{Ric}=\tfrac{k+l-2}2g,\qquad
    \langle R(X,Y)Y,X\rangle=|[X,Y]|^2\ge0.
  \]
  The isotropy representation on $t_p$ is irreducible but is transitive on its
  unit sphere only when $k=1$ or $l=1$; precisely then the curvature is positive
  and constant. Moreover
  \cref{cor:pet-ch10-maximal-subalgebra} gives
  $\mathrm{iso}=\mathfrak{so}(k+l)$.
\end{example}
\begin{proof}
  \uses{ex:pet-ch10-compact-grassmannian}
  \sketch
  Transitivity of $O(k+l)$ and the graph chart at $\mathbb R^k$ give the stated
  quotient, tangent space, and isotropy action. The element $(I_k,-I_l)$ acts
  by $-I$ on graph maps. Block decomposition then identifies $t_p$ with the
  off-diagonal skew matrices. Since
  $B(X,Y)=(k+l-2)\operatorname{tr}(XY)$,
  \cref{thm:pet-ch10-curvature-bracket-formula} yields
  \[
    \operatorname{Ric}=\frac{k+l-2}{2}g,\qquad
    \langle R(X,Y)Y,X\rangle=|[X,Y]|^2.
  \]
  Irreducibility of the isotropy module and
  \cref{cor:pet-ch10-maximal-subalgebra} identify the isometry algebra with
  $\mathfrak{so}(k+l)$.
\end{proof}

\begin{example}[the hyperbolic Grassmannian]
  \label{ex:pet-ch10-hyperbolic-grassmannian}
  \lean{PetersenLib.hyperbolicGrassmannianSymmetricSpace}
  \source{petersen:page-0392,petersen:page-0393}
  \uses{ex:pet-ch10-compact-grassmannian}
  \dcref{ch10:10.2.2}
  On $\mathbb R^{k,l}$ use
  $v^\top I_{k,l}w=\sum_{i\le k}v_iw_i-\sum_{i>k}v_iw_i$. Then
  \[
    O(k,l)=\{X:XI_{k,l}X^\top=I_{k,l}\},
  \]
  and
  \[
    \mathfrak{so}(k,l)
    =\left\{\begin{pmatrix}Y_1&B\\B^\top&Y_2\end{pmatrix}:
      Y_1\in\mathfrak{so}(k),\ Y_2\in\mathfrak{so}(l)\right\}.
  \]
  The oriented positive $k$-planes form the symmetric space
  \[
    \widetilde G_k(\mathbb R^{k,l})
      =O(k,l)/(\mathrm{SO}(k)\times O(l)).
  \]
  At $p=\mathbb R^k$,
  $t_p=\{\left(\begin{smallmatrix}0&B\\B^\top&0\end{smallmatrix}\right)\}$ and
  $g(X,Y)=\operatorname{tr}(XY)=(k+l-2)^{-1}B(X,Y)$. Therefore
  \[
    \operatorname{Ric}=-\tfrac{k+l-2}2g,\qquad
    \langle R(X,Y)Y,X\rangle=-|[X,Y]|^2\le0,
  \]
  and $\mathrm{iso}=\mathfrak{so}(k,l)$. The case $l=1$ is hyperbolic space
  $H^k$.
\end{example}
\begin{proof}
  \uses{ex:pet-ch10-hyperbolic-grassmannian}
  \sketch
  Repeat the quotient and involution argument of
  \cref{ex:pet-ch10-compact-grassmannian} with $O(k,l)$. The off-diagonal
  matrices are now symmetric, so
  $g=(k+l-2)^{-1}B$ on $t_p$. Hence
  \cref{thm:pet-ch10-curvature-bracket-formula} reverses the compact signs and
  gives the displayed Ricci and sectional-curvature formulas. The same
  isotropy irreducibility, together with
  \cref{cor:pet-ch10-maximal-subalgebra}, gives
  $\mathrm{iso}=\mathfrak{so}(k,l)$.
\end{proof}

\begin{example}[complex projective space as a symmetric space]
  \label{ex:pet-ch10-complex-projective-space}
  \lean{PetersenLib.complexProjectiveSpaceSymmetricSpace}
  \source{petersen:page-0393,petersen:page-0394,petersen:page-0395}
  \uses{def:pet-ch10-tp-decomposition, thm:pet-ch10-curvature-bracket-formula}
  \dcref{ch10:10.2.3}
  For the first coordinate line $p\in\mathbb{CP}^n$,
  \[
    \mathbb{CP}^n
      =\mathrm{SU}(n+1)/\mathrm S(\mathrm U(1)\times\mathrm U(n)),
  \]
  where the isotropy is identified with $\mathrm U(n)$ by
  $A\mapsto\left(\begin{smallmatrix}\det A^{-1}&0\\0&A\end{smallmatrix}\right)$.
  The element
  $\left(\begin{smallmatrix}(-1)^n&0\\0&-I_n\end{smallmatrix}\right)$ induces
  the point symmetry. Moreover
  \[
    \mathfrak{su}(n+1)=\mathfrak t_p\oplus\mathfrak u(n),\qquad
    \mathfrak t_p=\Big\{\begin{pmatrix}0&-z^*\\z&0\end{pmatrix}\mid z\in\mathbb C^n\Big\},
  \]
  and $g(X,Y)=-\operatorname{tr}(XY)=-[2(n+1)]^{-1}B(X,Y)$ on $\mathfrak t_p$.
  Thus
  \[
    \operatorname{Ric}=(n+1)g,\qquad
    \langle R(X,Y)Y,X\rangle=|[X,Y]|^2=4|\operatorname{Im}(z^*w)|^2-2\operatorname{Re}(z^*w)^2+2|z|^2|w|^2.
  \]
  For an orthonormal plane, $|z|^2=|w|^2=\tfrac12$ and
  $\operatorname{Re}(z^*w)=0$, so the sectional curvature is
  $6|z^*w|^2+\tfrac12$. Hence the Fubini--Study metric satisfies
  $\tfrac14\le\sec\le2$, with its maximum at $w=iz$ and its minimum when
  $z^*w=0$.
\end{example}
\begin{proof}
  \uses{ex:pet-ch10-complex-projective-space}
  \sketch
  The unitary action on complex lines is transitive; quotienting its scalar
  kernel gives the stated special-unitary action and block-diagonal isotropy.
  The displayed order-two element negates the complementary tangent block.
  The Lie-algebra block form gives $t_p$ and
  $g=-[2(n+1)]^{-1}B$, so
  \cref{thm:pet-ch10-curvature-bracket-formula} gives
  $\operatorname{Ric}=(n+1)g$. For tangent vectors represented by $z,w$,
  \[
    [X,Y]=\begin{pmatrix}-(z^*w-w^*z)&0\\0&wz^*-zw^*\end{pmatrix}.
  \]
  Expanding its norm yields the stated curvature expression. Under the
  orthonormality conditions it becomes $6|z^*w|^2+\tfrac12$, whose extrema
  occur at $w=iz$ and $z^*w=0$.
\end{proof}

\begin{example}[$\mathrm{SL}(n)/\mathrm{SO}(n)$]
  \label{ex:pet-ch10-sl-so-quotient}
  \lean{PetersenLib.slSoSymmetricSpace}
  \source{petersen:page-0395,petersen:page-0396}
  \uses{def:pet-ch10-tp-decomposition, thm:pet-ch10-curvature-bracket-formula}
  \dcref{ch10:10.2.4}
  Decompose $\mathfrak{sl}(n)=\mathfrak t\oplus\mathfrak{so}(n)$ into symmetric
  traceless and skew-symmetric matrices. The involution
  $\sigma(X)=-X^\top$ negates $\mathfrak t$ and fixes $\mathfrak{so}(n)$. On
  $\mathfrak t$,
  \[
    g(X,Y)=\operatorname{tr}(XY)=\tfrac1{2n}B(X,Y),\qquad
    \operatorname{Ric}=-ng,
  \]
  and $\langle R(X,Y)Z,W\rangle=\langle[X,Y],[Z,W]\rangle$. Thus
  $\mathrm{SL}(n)/\mathrm{SO}(n)$ has nonpositive sectional curvature.
\end{example}
\begin{proof}
  \uses{ex:pet-ch10-sl-so-quotient}
  \sketch
  Transpose decomposes $\mathfrak{sl}(n)$ into symmetric traceless and
  skew-symmetric matrices, while $X\mapsto-X^\top$ is an involution with these
  as its $(-1)$- and $1$-eigenspaces. On the first summand,
  $g=(2n)^{-1}B$. Thus
  \cref{thm:pet-ch10-curvature-bracket-formula} gives
  $\operatorname{Ric}=-ng$ and
  \[
    \langle R(X,Y)Y,X\rangle=-\langle[X,Y],[X,Y]\rangle\le0.
  \]
  This is the noncompact-sign counterpart of
  \cref{ex:pet-ch10-compact-grassmannian}.
\end{proof}

\begin{example}[Lie groups as symmetric spaces]
  \label{ex:pet-ch10-lie-group-symmetric}
  \lean{PetersenLib.compactLieGroupSymmetricSpace,PetersenLib.noncompactDualSymmetricSpace}
  \source{petersen:page-0389,petersen:page-0396}
  \uses{rem:pet-ch10-symmetric-space-from-involution, def:pet-ch10-killing-form, thm:pet-ch10-curvature-bracket-formula}
  \dcref{ch10:10.2.5}
  Let $G$ be compact with a biinvariant metric and centerless Lie algebra
  $\mathfrak g$. Then $B$ is negative definite and $g=-B$ is biinvariant. The
  swap involution on $\mathfrak g\oplus\mathfrak g$ has eigenspaces
  \[
    \mathfrak k=\mathfrak g^\Delta,
    \qquad \mathfrak t=\{(X,-X):X\in\mathfrak g\},
  \]
  and realizes $G=(G\times G)/G^\Delta$. On $\mathfrak t$ the metric has the
  factor-two scaling, $\operatorname{Ric}=\tfrac12g$, and the curvature operator
  is nonnegative.

  The noncompact dual uses complex conjugation on
  $\mathfrak g\otimes\mathbb C=\mathfrak g\oplus i\mathfrak g$. Here
  $\mathfrak k=\mathfrak g$, $\mathfrak t=i\mathfrak g$,
  \[
    g(iX,iY)=-B(X,Y)=B(iX,iY),\qquad
    \operatorname{Ric}=-\tfrac12g,
  \]
  and the curvature operator is nonpositive.
\end{example}
\begin{proof}
  \uses{ex:pet-ch10-lie-group-symmetric}
  \sketch
  For a biinvariant metric, $\mathrm{ad}_X=2(\nabla X)_e$ is skew-adjoint.
  Hence $B\le0$, with kernel the center, as in
  \cref{rem:pet-ch10-adjoint-skew-symmetric}. The swap involution on
  $\mathfrak g\oplus\mathfrak g$ has diagonal and antidiagonal eigenspaces, so
  \cref{rem:pet-ch10-symmetric-space-from-involution} realizes
  $G=(G\times G)/G^\Delta$. Applying
  \cref{thm:pet-ch10-curvature-bracket-formula} with the doubled Killing form
  gives the compact curvature sign and Ricci constant.

  Complex conjugation on $\mathfrak g\otimes\mathbb C$ fixes $\mathfrak g$ and
  negates $i\mathfrak g$. Since $B(iX,iY)=-B(X,Y)$, it defines the positive
  metric on the noncompact tangent summand, and the same bracket formula gives
  $\operatorname{Ric}=-\tfrac12g$ and nonpositive curvature.
\end{proof}

\section{Holonomy}

\subsection{The Holonomy Group}

\begin{definition}[holonomy group and restricted holonomy group]
  \label{def:pet-ch10-holonomy-group}
  \lean{PetersenLib.HolonomyGroup,PetersenLib.RestrictedHolonomyGroup}
  \source{petersen:page-0396,petersen:page-0397}
  Let $P_{c(a)}^{c(b)}$ denote parallel translation along a piecewise smooth
  curve $c:[a,b]\to M$. For fixed $p\in M$, the maps $P_p^p$ arising from loops
  based at $p$ form the holonomy group
  $\mathrm{Hol}_p(M,g)\subset O(T_pM)$. The subgroup obtained from contractible
  loops is the restricted holonomy group $\mathrm{Hol}_p^0(M,g)$; it is a
  connected compact normal Lie subgroup. The full holonomy group is also a Lie
  group and may be noncompact.
\end{definition}

\begin{remark}[elementary properties of the holonomy group]
  \label{rem:pet-ch10-holonomy-properties}
  \lean{PetersenLib.holonomyEuclidean,PetersenLib.holonomySphere,PetersenLib.holonomyHyperbolicSpace,PetersenLib.holonomyOrientability,PetersenLib.holonomyUniversalCover,PetersenLib.holonomyProduct,PetersenLib.holonomyConjugate,PetersenLib.holonomyInvariantTensor}
  \source{petersen:page-0397}
  \uses{def:pet-ch10-holonomy-group}
  \begin{enumerate}
    \item[(a)] $\mathrm{Hol}_p(\mathbb R^n)=\{I\}$.
    \item[(b)] $\mathrm{Hol}_p(S^n(R))=\mathrm{SO}(n)$.
    \item[(c)] $\mathrm{Hol}_p(H^n)=\mathrm{SO}(n)$.
    \item[(d)] $\mathrm{Hol}_p(M,g)\subset\mathrm{SO}(n)$ iff $M$ is orientable.
    \item[(e)] $\mathrm{Hol}_p(\widetilde M,\widetilde g)=\mathrm{Hol}_p^0(M,g)$ for the
      universal cover $\widetilde M$.
    \item[(f)] $\mathrm{Hol}_{(p,q)}(M_1\times M_2,g_1+g_2)=\mathrm{Hol}_p(M_1,g_1)\times\mathrm{Hol}_q(M_2,g_2)$.
    \item[(g)] $\mathrm{Hol}_p(M,g)$ is conjugate to $\mathrm{Hol}_q(M,g)$ via parallel
      translation along any curve from $p$ to $q$.
    \item[(h)] A tensor at $p$ extends to a parallel tensor on $(M,g)$ iff it is
      $\mathrm{Hol}_p(M,g)$-invariant; e.g.\ a $2$-form $\omega$ extends iff
      $\omega(Pv,Pw)=\omega(v,w)$ for all $P\in\mathrm{Hol}_p(M,g)$,
      $v,w\in T_pM$.
  \end{enumerate}
\end{remark}

An orthogonal decomposition of $T_pM$ into irreducible
$\mathrm{Hol}_p^0$-modules transports to a parallel splitting
$TM=\eta_1\oplus\cdots\oplus\eta_k$. Product metrics give the corresponding
product decomposition of holonomy.

\begin{theorem}[de Rham decomposition]
  \label{thm:pet-ch10-de-rham-decomposition}
  \lean{PetersenLib.deRhamDecomposition}
  \source{petersen:page-0397,petersen:page-0398}
  \uses{rem:pet-ch10-holonomy-properties}
  \dcref{ch10:10.3.1}
  Let $(M,g)$ be a Riemannian manifold and $TM=\eta_1\oplus\cdots\oplus\eta_k$
  the decomposition of the tangent bundle into irreducible components for the
  restricted holonomy. Around each $p\in M$ there is a neighborhood $U$ with a
  product structure $(U,g)=(U_1\times\cdots\times U_k,g_1+\cdots+g_k)$,
  $TU_i=\eta_i|_U$. If $(M,g)$ is simply connected and complete, this splitting
  is global: $(M,g)=(M_1\times\cdots\times M_k,g_1+\cdots+g_k)$, $TM_i=\eta_i$.
\end{theorem}
\begin{proof}
  \uses{thm:pet-ch10-de-rham-decomposition}
  \sketch
  Parallel invariance makes each $\eta_i$ closed under the connection. Since
  the Levi--Civita connection is torsion-free, Frobenius gives integral leaves;
  connection invariance also makes those leaves totally geodesic. Orthogonality
  and parallelism of the summands then split the metric on a product
  neighborhood.

  Through a fixed point, take the maximal leaf $M_i$ of each $\eta_i$. These
  leaves are complete when $M$ is complete. The local product isometry extends
  to a Riemannian covering between $M$ and the product of the $M_i$; completeness
  provides global continuation. If $M$ is simply connected, passage to the
  universal covers of the factors turns this covering into the asserted global
  isometry.
\end{proof}

\subsection{Rough Classification of Symmetric Spaces}

A Riemannian manifold is \emph{irreducible} when its holonomy representation
has no proper invariant subspace. The de Rham theorem reduces classification
questions to this case.

\begin{theorem}[parallel tensors on irreducible manifolds]
  \label{thm:pet-ch10-parallel-tensor-eigenvalue}
  \lean{PetersenLib.parallelTensorSingleEigenvalue}
  \source{petersen:page-0399}
  \uses{thm:pet-ch10-de-rham-decomposition}
  \dcref{ch10:10.3.2}
  Let $(M,g)$ be an irreducible Riemannian manifold with a parallel $(1,1)$-tensor
  $T$. Then the symmetric part $S=\tfrac12(T+T^*)$ and skew-symmetric part
  $A=\tfrac12(T-T^*)$ each have exactly one (complex) eigenvalue. If $A\ne0$, it
  induces (after scaling) a parallel complex structure, i.e.\ a Kähler structure.
\end{theorem}
\begin{proof}
  \uses{thm:pet-ch10-parallel-tensor-eigenvalue}
  \sketch
  Metric compatibility implies
  \[
    g((\nabla_XT)Y,Z)=g(Y,(\nabla_XT^*)Z),
  \]
  so $T^*$, $S$, and $A$ are parallel. Every real eigenspace of $S$, and every
  real invariant subspace associated with a conjugate pair of eigenvalues of
  $A$, is therefore preserved by holonomy. Irreducibility leaves only one such
  summand in each decomposition. Thus $S$ is scalar and the nonzero
  skew-adjoint $A$ satisfies $A^2=-\mu^2I$ for one $\mu>0$. Then
  $J=\mu^{-1}A$ is parallel and $J^2=-I$, which is the required Kähler
  structure.
\end{proof}

\begin{corollary}[irreducible symmetric spaces are Einstein]
  \label{cor:pet-ch10-irreducible-symmetric-einstein}
  \lean{PetersenLib.irreducibleSymmetricSpaceEinstein}
  \source{petersen:page-0399}
  \uses{thm:pet-ch10-parallel-tensor-eigenvalue}
  \dcref{ch10:10.3.3}
  A simply connected irreducible symmetric space is an Einstein manifold; in
  particular its curvature operator is nonnegative or nonpositive according to
  the sign of the Einstein constant.
\end{corollary}

\begin{remark}[compact, flat, and noncompact type]
  \label{rem:pet-ch10-compact-flat-noncompact-type}
  \lean{PetersenLib.symmetricSpaceCompactFlatNoncompactType}
  \source{petersen:page-0399,petersen:page-0400}
  \uses{cor:pet-ch10-irreducible-symmetric-einstein}
  For an irreducible symmetric space, \cref{cor:pet-ch10-irreducible-symmetric-einstein}
  gives an Einstein constant, and three cases occur. \emph{Compact type}: if the
  Einstein constant is positive, Myers' diameter bound forces $M$ compact, with
  nonnegative curvature operator. \emph{Flat type}: if $M$ is Ricci flat, it is
  flat, and the only irreducible Ricci-flat examples are $S^1$ and $\mathbb R^1$.
  \emph{Noncompact type}: if the Einstein constant is negative, Bochner's theorem
  on Killing fields forces $M$ noncompact, with nonpositive curvature operator.
\end{remark}

\begin{definition}[Types I--IV of irreducible symmetric spaces]
  \label{def:pet-ch10-symmetric-space-type-I-IV}
  \lean{PetersenLib.symmetricSpaceTypeI,PetersenLib.symmetricSpaceTypeII,PetersenLib.symmetricSpaceTypeIII,PetersenLib.symmetricSpaceTypeIV}
  \source{petersen:page-0400}
  \uses{rem:pet-ch10-compact-flat-noncompact-type, rem:pet-ch10-symmetric-space-from-involution}
  Irreducible symmetric spaces of compact and noncompact type occur in dual
  pairs described by Lie algebra pairs $(\mathfrak g,\mathfrak t_p)$ with an
  involution having $1$-eigenspace $\mathfrak k$ (the Lie algebra of a compact
  group acting on the $(-1)$-eigenspace $\mathfrak t_p$):
  \begin{itemize}
    \item \textbf{Type I}: compact, $(\mathfrak g,\mathfrak t)$ with $\mathfrak g$
      simple, the Lie algebra of a compact group, $\mathfrak t\subset\mathfrak g$
      a maximal subalgebra; e.g.\ $(\mathfrak{so}(k+l),\mathfrak{so}(k)\times\mathfrak{so}(l))$.
    \item \textbf{Type II}: compact, $(\mathfrak t\oplus\mathfrak t,\Delta\mathfrak t)$
      with $\mathfrak t$ simple, the Lie algebra of a compact Lie group $G$; the
      space is $G$ with a biinvariant metric.
    \item \textbf{Type III}: noncompact, $(\mathfrak g,\mathfrak t)$ with
      $\mathfrak g$ simple, the Lie algebra of a noncompact group,
      $\mathfrak t\subset\mathfrak g$ maximal corresponding to a compact group;
      e.g.\ $(\mathfrak{so}(k,l),\mathfrak{so}(k)\times\mathfrak{so}(l))$ or
      $(\mathfrak{sl}(n),\mathfrak{so}(n))$.
    \item \textbf{Type IV}: noncompact, $(\mathfrak t\otimes\mathbb C,\mathfrak t)$
      with $\mathfrak t$ simple, corresponding to a compact group, and
      $\mathfrak t\otimes\mathbb C=\mathfrak t\oplus i\mathfrak t$ its
      complexification; e.g.\ $(\mathfrak{so}(n,\mathbb C),\mathfrak{so}(n))$.
  \end{itemize}
  In every case \cref{cor:pet-ch10-maximal-subalgebra} shows $t_p=\mathrm{iso}_p$.
  Since compact-type symmetric spaces have nonnegative curvature operator, the
  Bochner technique shows all their harmonic forms are parallel, hence
  holonomy-invariant, reducing the computation of their cohomology to a
  classical invariance problem for the isotropy representation — used, e.g., to
  compute Pontryagin and Chern classes from the cohomology of real and complex
  Grassmannians.
\end{definition}

\subsection{Curvature and Holonomy}

\begin{definition}[Berger's list of irreducible holonomy groups]
  \label{def:pet-ch10-berger-holonomy-table}
  \lean{PetersenLib.bergerHolonomyList}
  \source{petersen:page-0401}
  For a simply connected irreducible Riemannian $n$-manifold whose holonomy
  acts transitively on the unit sphere of $T_pM$, the possible holonomy groups
  $\mathrm{Hol}_p$, and the resulting geometric structure, are:
  \begin{center}
    \begin{tabular}{c|c|c}
      $\dim=n$ & $\mathrm{Hol}_p$ & properties \\\hline
      $n$ & $\mathrm{SO}(n)$ & generic case \\
      $n=2m$ & $\mathrm U(m)$ & Kähler \\
      $n=2m$ & $\mathrm{SU}(m)$ & Kähler and Ricci flat \\
      $n=4m$ & $\mathrm{Sp}(1)\cdot\mathrm{Sp}(m)$ & quaternionic-Kähler and Einstein \\
      $n=4m$ & $\mathrm{Sp}(m)$ & hyper-Kähler and Ricci flat \\
      $n=16$ & $\mathrm{Spin}(9)$ & symmetric and Einstein \\
      $n=8$ & $\mathrm{Spin}(7)$ & Ricci flat \\
      $n=7$ & $G_2$ & Ricci flat
    \end{tabular}
  \end{center}
\end{definition}

\begin{theorem}[holonomy classification alternative]
  \label{thm:pet-ch10-berger-simons-classification}
  \lean{PetersenLib.bergerSimonsClassification}
  \source{petersen:page-0401}
  \uses{def:pet-ch10-berger-holonomy-table, def:pet-ch10-rank}
  \dcref{ch10:10.3.4}
  Let $(M,g)$ be a simply connected irreducible Riemannian $n$-manifold. Then
  $\mathrm{Hol}_p$ either acts transitively on the unit sphere of $T_pM$, in
  which case $\mathrm{Hol}_p$ is one of the groups of
  \cref{def:pet-ch10-berger-holonomy-table}, or $(M,g)$ is a symmetric space of
  rank $\ge2$.
\end{theorem}

\begin{corollary}[nontransitive holonomy forces reducibility or rank $\ge2$]
  \label{cor:pet-ch10-nontransitive-holonomy}
  \lean{PetersenLib.nontransitiveHolonomyReducibleOrSymmetric}
  \source{petersen:page-0402}
  \uses{thm:pet-ch10-berger-simons-classification}
  \dcref{ch10:10.3.5}
  If a Riemannian manifold's holonomy does not act transitively on the unit
  sphere, then it is either reducible or a locally symmetric space of rank
  $\ge2$; in particular its rank is $\ge2$.
\end{corollary}

\begin{theorem}[rank rigidity in nonpositive curvature]
  \label{thm:pet-ch10-rank-rigidity-nonpositive}
  \lean{PetersenLib.rankRigidityNonpositiveCurvature}
  \source{petersen:page-0402}
  \uses{cor:pet-ch10-nontransitive-holonomy}
  \dcref{ch10:10.3.6}
  A compact Riemannian manifold of nonpositive curvature and rank $\ge2$ does
  not have transitive holonomy; in particular it is either reducible or locally
  symmetric.
\end{theorem}

\begin{theorem}[classification under nonnegative curvature operator]
  \label{thm:pet-ch10-gallot-meyer}
  \lean{PetersenLib.gallotMeyerNonnegativeCurvatureOperator}
  \source{petersen:page-0402,petersen:page-0403}
  \uses{thm:pet-ch10-berger-simons-classification, thm:pet-ch10-parallel-tensor-eigenvalue, thm:pet-ch10-de-rham-decomposition}
  \dcref{ch10:10.3.7}
  Let $(M,g)$ be a compact Riemannian $n$-manifold with nonnegative curvature
  operator. Then one of:
  \begin{enumerate}
    \item[(a)] $(M,g)$ is reducible or locally symmetric;
    \item[(b)] $\mathrm{Hol}^0(M,g)=\mathrm{SO}(n)$ and the universal cover is a
      homology sphere;
    \item[(c)] $\mathrm{Hol}^0(M,g)=\mathrm U(n/2)$ and the universal cover is a
      homology $\mathbb{CP}^{n/2}$.
  \end{enumerate}
\end{theorem}
\begin{proof}
  \uses{thm:pet-ch10-gallot-meyer}
  \sketch
  If $\pi_1(M)$ is infinite, the splitting theorem for nonnegative Ricci
  curvature gives an Euclidean factor in the universal cover, hence case (a).
  Otherwise pass to the simply connected cover. Bochner turns every harmonic
  form into a parallel form.

  Assume irreducibility. By
  \cref{thm:pet-ch10-berger-simons-classification}, nontransitive holonomy gives
  local symmetry. Among the transitive groups in
  \cref{def:pet-ch10-berger-holonomy-table}, Ricci-flat cases are flat, while
  the quaternionic and $\mathrm{Spin}(9)$ cases are Einstein and hence locally
  symmetric under the curvature-operator hypothesis. Thus only
  $\mathrm{SO}(n)$ and $\mathrm U(n/2)$ remain.

  The group $\mathrm{SO}(n)$ has no invariant $p$-forms for $0<p<n$: reverse
  one argument vector and compensate the determinant in an unused direction.
  Hence there is no intermediate cohomology. In the unitary case the Kähler
  form is parallel; invariant-form theory, with the uniqueness mechanism of
  \cref{thm:pet-ch10-parallel-tensor-eigenvalue}, shows that odd cohomology
  vanishes and each $H^{2k}$ is generated by the nonzero power $\omega^k$.
  These are precisely the cohomology groups of $\mathbb{CP}^{n/2}$.
\end{proof}

\begin{theorem}[classification under nonnegative complex sectional curvature]
  \label{thm:pet-ch10-brendle-schoen}
  \lean{PetersenLib.brendleSchoenNonnegativeComplexSectionalCurvature}
  \source{petersen:page-0404}
  \uses{thm:pet-ch10-gallot-meyer}
  \dcref{ch10:10.3.8}
  Let $(M,g)$ be a compact Riemannian $n$-manifold with nonnegative complex
  sectional curvature. Then one of:
  \begin{enumerate}
    \item[(a)] $(M,g)$ is reducible or locally symmetric;
    \item[(b)] $M$ is diffeomorphic to a space of constant positive curvature;
    \item[(c)] the universal cover of $M$ is biholomorphic to $\mathbb{CP}^{n/2}$.
  \end{enumerate}
\end{theorem}

\section{Exercises}

\begin{remark}[geodesic loops are closed geodesics]
  \label{rem:pet-ch10-exercise-symmetric-flow-axis}
  \lean{PetersenLib.exercise10_5_1}
  \source{petersen:page-0405}
  \uses{def:pet-ch10-tp-decomposition}
  \dcref{ch10:10.5.1}
  Let $M$ be a symmetric space and $X\in t_p$ nontrivial (so $X$ is a Killing
  field with $(\nabla X)|_p=0$). Show: (1) the flow of $X$ is
  $F_s=A_{\exp(sX|_p)}\circ A_p$; (2) $c(t)=\exp(tX|_p)$ is an axis for $F_s$,
  i.e.\ $F_s(c(t))=c(t+s)$; (3) if $c(a)=c(b)$ then $\dot c(a)=\dot c(b)$;
  (4) conclude that geodesic loops are always closed geodesics.
\end{remark}

\begin{remark}[the fundamental group of a symmetric space is abelian]
  \label{rem:pet-ch10-exercise-pi1-abelian}
  \lean{PetersenLib.exercise10_5_2}
  \source{petersen:page-0405}
  \uses{rem:pet-ch10-exercise-symmetric-flow-axis}
  \dcref{ch10:10.5.2}
  Let $M$ be a symmetric space. Show that $A_p$ acting on $\pi_1(M,p)$ by
  $[\gamma]\mapsto[\gamma]^{-1}$ shows $\pi_1(M,p)$ is abelian. (Hint: every
  element of $\pi_1(M,p)$ is represented by a geodesic loop, which by
  \cref{rem:pet-ch10-exercise-symmetric-flow-axis} is a closed geodesic.)
\end{remark}

\begin{remark}[Jacobi fields along a geodesic in a symmetric space]
  \label{rem:pet-ch10-exercise-jacobi-fields-eigenvector}
  \lean{PetersenLib.exercise10_5_3}
  \source{petersen:page-0405,petersen:page-0406}
  \uses{def:pet-ch10-tp-decomposition}
  \dcref{ch10:10.5.3}
  Let $M$ be a symmetric space and $c$ a geodesic. (1) If $E(t)$ is parallel
  along $c$, show $R(E,\dot c)\dot c$ is parallel. (2) If $E(0)$ is an eigenvector
  of $R(\cdot,\dot c(0))\dot c(0)$ with eigenvalue $\kappa$, show
  $J_1(t)=\operatorname{sn}_\kappa(t)E(t)$ and $J_2(t)=\operatorname{sn}_\kappa'(t)E(t)$ are Jacobi
  fields along $c$. (3) If $J$ is a Jacobi field with $J(0)=0$, $J'(t_0)=0$, show
  $J(t_0/u)=0$ for the appropriate rescaling parameter $u$. (4) With $J$ as in
  (3), construct a geodesic variation $c(s,t)$ with $\partial_sc(0,t)=J(t)$,
  $c(s,0)=c(0)$, $c(s,t_0)=c(t_0)$.
\end{remark}

\begin{remark}[nonnegative curvature of compact symmetric spaces]
  \label{rem:pet-ch10-exercise-compact-nonneg-curvature}
  \lean{PetersenLib.exercise10_5_4}
  \source{petersen:page-0406}
  \uses{rem:pet-ch10-exercise-jacobi-fields-eigenvector}
  \dcref{ch10:10.5.4}
  Let $M$ be a symmetric space. (1) Show directly that if $M$ is compact then
  $\sec\ge0$ (hint: argue by contradiction, producing an unbounded Jacobi field
  along a geodesic). (2) Show that if $c$ is a closed geodesic then
  $R(\cdot,\dot c)\dot c$ has no negative eigenvalues. (3) Show that if
  $\operatorname{Ric}>0$ then $\sec\ge0$.
\end{remark}

\begin{remark}[equivalent characterizations of a flat parallel direction]
  \label{rem:pet-ch10-exercise-jacobi-field-equivalence}
  \lean{PetersenLib.exercise10_5_5}
  \source{petersen:page-0405,petersen:page-0406}
  \uses{def:pet-ch10-holonomy-group}
  \dcref{ch10:10.5.5}
  Assume $M$ has nonpositive or nonnegative sectional curvature, $c$ a geodesic,
  $E$ a parallel field along $c$. Show the following are equivalent: (1)
  $g(R(E,\dot c)\dot c,E)=0$ everywhere along $c$; (2) $R(E,\dot c)\dot c=0$
  everywhere along $c$; (3) $E$ is a Jacobi field.
\end{remark}

\begin{remark}[Killing forms of the classical Lie algebras]
  \label{rem:pet-ch10-exercise-killing-form-classical}
  \lean{PetersenLib.exercise10_5_6}
  \source{petersen:page-0406}
  \uses{def:pet-ch10-killing-form}
  \dcref{ch10:10.5.6}
  (1) Show the Killing form of $\mathfrak g\otimes\mathbb C$ is the
  complexification of the Killing form $B$ of a real Lie algebra $\mathfrak g$.
  (2) Show the Killing forms of $\mathfrak{gl}(n,\mathbb C)$ and $\mathfrak{gl}(n)$
  are $B(X,Y)=2n\operatorname{tr}(XY)-2\operatorname{tr}X\operatorname{tr}Y$ (hint: use the basis
  $E_{ij}=[\delta_{i\delta_j}]$). (3) Show that on
  $\mathfrak{sl}(n,\mathbb C)=\mathfrak{sl}(n)\otimes\mathbb C$, $B(X,Y)=2n\operatorname{tr}(XY)$
  (hint: use (2) and that $I\in\mathfrak{gl}(n,\mathbb C)$ commutes with
  $\mathfrak{sl}(n,\mathbb C)$). (4) Show $\mathfrak{sl}(k+l,\mathbb C)=\mathfrak{su}(k,l)\otimes\mathbb C$,
  and conclude $\mathfrak{sl}(k+l)$ and $\mathfrak{su}(k,l)$ have Killing form
  $B(X,Y)=2(k+l)\operatorname{tr}(XY)$. (5) Show that on
  $\mathfrak{so}(n,\mathbb C)=\mathfrak{so}(n)\otimes\mathbb C$, $B(X,Y)=(n-2)\operatorname{tr}(XY)$
  (hint: use the basis $E_{ij}-E_{ji}$, $i<j$). (6) Show
  $\mathfrak{so}(k+l,\mathbb C)=\mathfrak{so}(k,l)\otimes\mathbb C$, and conclude
  $\mathfrak{so}(k,l)$ has Killing form $B(X,Y)=(k+l-2)\operatorname{tr}(XY)$.
\end{remark}

\begin{remark}[nondegenerate bilinear forms as a symmetric space]
  \label{rem:pet-ch10-exercise-gl-so-symmetric}
  \lean{PetersenLib.exercise10_5_7}
  \source{petersen:page-0406}
  \uses{rem:pet-ch10-symmetric-space-from-involution}
  \dcref{ch10:10.5.7}
  Show that $\mathrm{GL}^+(p+q,\mathbb R)/\mathrm{SO}(p,q)$ defines a symmetric space,
  identifiable with the nondegenerate bilinear forms on $\mathbb R^{p+q}$ of index
  $q$.
\end{remark}

\begin{remark}[$\mathrm U(p,q)/\mathrm{SO}(p,q)$ as a symmetric space]
  \label{rem:pet-ch10-exercise-u-pq-so-pq-symmetric}
  \lean{PetersenLib.exercise10_5_8}
  \source{petersen:page-0406}
  \uses{rem:pet-ch10-symmetric-space-from-involution}
  \dcref{ch10:10.5.8}
  Show that $\mathrm U(p,q)/\mathrm{SO}(p,q)$ defines a symmetric space and that
  $\mathfrak u(p,q)\otimes\mathbb C=\mathfrak{gl}(p+q,\mathbb C)$.
\end{remark}

\begin{remark}[holonomy of $\mathbb{CP}^n$]
  \label{rem:pet-ch10-exercise-cpn-holonomy}
  \lean{PetersenLib.exercise10_5_9}
  \source{petersen:page-0406}
  \uses{ex:pet-ch10-complex-projective-space, def:pet-ch10-holonomy-group}
  \dcref{ch10:10.5.9}
  Show that the holonomy of $\mathbb{CP}^n$ is $\mathrm U(n)$.
\end{remark}

\begin{remark}[covering spaces of symmetric spaces]
  \label{rem:pet-ch10-exercise-covering-symmetric}
  \lean{PetersenLib.exercise10_5_10}
  \source{petersen:page-0406}
  \uses{prop:pet-ch10-sigma-p-involution}
  \dcref{ch10:10.5.10}
  Show that a covering space of a symmetric space is a symmetric space; give an
  example showing the converse fails.
\end{remark}

\begin{remark}[flatness iff discrete holonomy]
  \label{rem:pet-ch10-exercise-flat-discrete-holonomy}
  \lean{PetersenLib.exercise10_5_11}
  \source{petersen:page-0406}
  \uses{def:pet-ch10-holonomy-group}
  \dcref{ch10:10.5.11}
  Show that a manifold is flat if and only if the (restricted) holonomy Lie
  algebra $\mathfrak{hol}_p=\{0\}$, i.e.\ the holonomy group is discrete.
\end{remark}

\begin{remark}[finite fundamental group under irreducible holonomy and $\operatorname{Ric}\ge0$]
  \label{rem:pet-ch10-exercise-finite-fundamental-group}
  \lean{PetersenLib.exercise10_5_12}
  \source{petersen:page-0406}
  \uses{def:pet-ch10-holonomy-group}
  \dcref{ch10:10.5.12}
  Show that a compact Riemannian manifold with irreducible restricted holonomy
  and $\operatorname{Ric}\ge0$ has finite fundamental group.
\end{remark}

\begin{remark}[identifying $\mathrm{SL}(2,\mathbb R)/\mathrm{SO}(2)$ and $\mathrm{SL}(2,\mathbb C)/\mathrm{SU}(2)$]
  \label{rem:pet-ch10-exercise-sl2-quotients}
  \lean{PetersenLib.exercise10_5_13}
  \source{petersen:page-0406}
  \uses{ex:pet-ch10-sl-so-quotient}
  \dcref{ch10:10.5.13}
  Which known spaces are described by $\mathrm{SL}(2,\mathbb R)/\mathrm{SO}(2)$ and
  $\mathrm{SL}(2,\mathbb C)/\mathrm{SU}(2)$?
\end{remark}

\begin{remark}[Kähler structure iff holonomy in $\mathrm U(m)$]
  \label{rem:pet-ch10-exercise-kahler-unitary-holonomy}
  \lean{PetersenLib.exercise10_5_14}
  \source{petersen:page-0406}
  \uses{def:pet-ch10-holonomy-group, def:pet-ch10-berger-holonomy-table}
  \dcref{ch10:10.5.14}
  Show that the holonomy of a Riemannian manifold is contained in $\mathrm U(m)$
  if and only if it has a Kähler structure.
\end{remark}

\begin{remark}[maximal isotropy forces constant curvature]
  \label{rem:pet-ch10-exercise-isop-so-constant-curvature}
  \lean{PetersenLib.exercise10_5_15}
  \source{petersen:page-0406}
  \uses{def:pet-ch10-tp-decomposition}
  \dcref{ch10:10.5.15}
  Show that if a homogeneous space has $\mathrm{iso}_p=\mathfrak{so}(T_pM)$ at
  some $p$, then it has constant curvature.
\end{remark}

\begin{remark}[maximality of two subalgebras of $\mathfrak{so}(n)$]
  \label{rem:pet-ch10-exercise-maximal-subalgebras-son}
  \lean{PetersenLib.exercise10_5_16}
  \source{petersen:page-0406}
  \uses{ex:pet-ch10-compact-grassmannian, ex:pet-ch10-complex-projective-space}
  \dcref{ch10:10.5.16}
  Show that the subalgebras $\mathfrak{so}(k)\times\mathfrak{so}(n-k)$ and
  $\mathfrak u(n/2)$ are maximal in $\mathfrak{so}(n)$.
\end{remark}

\begin{remark}[$\mathfrak{su}(m)\subset\mathfrak{so}(m^2-1)$]
  \label{rem:pet-ch10-exercise-sum-in-so}
  \lean{PetersenLib.exercise10_5_17}
  \source{petersen:page-0406}
  \uses{def:pet-ch10-killing-form}
  \dcref{ch10:10.5.17}
  Show that $\mathfrak{su}(m)\subset\mathfrak{so}(m^2-1)$. (Hint: let
  $\mathfrak{su}(m)$ act on itself by the adjoint representation.)
\end{remark}

\begin{remark}[$t_p\subset\mathfrak{hol}_p$ in general]
  \label{rem:pet-ch10-exercise-tp-subset-holp}
  \lean{PetersenLib.exercise10_5_18}
  \source{petersen:page-0407}
  \uses{def:pet-ch10-tp-decomposition, def:pet-ch10-holonomy-group}
  \dcref{ch10:10.5.18}
  Show that for any Riemannian manifold $t_p\subset\mathfrak{hol}_p$ (the
  holonomy Lie algebra). Give an example where equality fails.
\end{remark}

\begin{remark}[$t_p=\mathfrak{hol}_p$ for symmetric spaces]
  \label{rem:pet-ch10-exercise-tp-equals-holp-symmetric}
  \lean{PetersenLib.exercise10_5_19}
  \source{petersen:page-0407}
  \uses{rem:pet-ch10-exercise-tp-subset-holp, cor:pet-ch10-maximal-subalgebra}
  \dcref{ch10:10.5.19}
  Show that for a symmetric space, $t_p=\mathfrak{hol}_p$. Use this to show
  that, unless the curvature is constant, $t_p=\mathfrak{hol}_p=\mathrm{iso}_p$
  provided $t_p\subset\mathfrak{so}(T_pM)$ is maximal.
\end{remark}

\begin{remark}[nonpositive curvature operator of two further quotients]
  \label{rem:pet-ch10-exercise-son-c-sln-c-symmetric}
  \lean{PetersenLib.exercise10_5_20}
  \source{petersen:page-0407}
  \uses{ex:pet-ch10-lie-group-symmetric}
  \dcref{ch10:10.5.20}
  Show that $\mathrm{SO}(n,\mathbb C)/\mathrm{SO}(n)$ and
  $\mathrm{SL}(n,\mathbb C)/\mathrm{SU}(n)$ are symmetric spaces with nonpositive
  curvature operator.
\end{remark}

\begin{remark}[quaternionic projective space as $\mathrm{Sp}(n+1)/(\mathrm{Sp}(1)\times\mathrm{Sp}(n))$]
  \label{rem:pet-ch10-exercise-quaternionic-projective-space}
  \lean{PetersenLib.exercise10_5_21}
  \source{petersen:page-0407}
  \uses{ex:pet-ch10-complex-projective-space}
  \dcref{ch10:10.5.21}
  The \emph{quaternionic projective space} $\mathbb{HP}^n$ is the space of
  quaternionic lines in $\mathbb H^{n+1}$. Define $\mathrm{Sp}(n)\subset\mathrm{SU}(2n)\subset\mathrm{SO}(4n)$
  as the orthogonal matrices commuting with the three complex structures
  generated by $i,j,k$ on $\mathbb R^{4n}$, equivalently the $n\times n$
  quaternionic matrices $A$ with $A^{-1}=A^*$. Show that, viewing $\mathbb H^{n+1}$
  as a right (or left) $\mathbb H$-module, $\mathbb{HP}^n=\mathrm{Sp}(n+1)/(\mathrm{Sp}(1)\times\mathrm{Sp}(n))$.
\end{remark}

\begin{remark}[hyperbolic analogues of complex projective space]
  \label{rem:pet-ch10-exercise-hyperbolic-complex-projective}
  \lean{PetersenLib.exercise10_5_22}
  \source{petersen:page-0407}
  \uses{ex:pet-ch10-complex-projective-space, ex:pet-ch10-hyperbolic-grassmannian}
  \dcref{ch10:10.5.22}
  Construct the hyperbolic analogues of the complex projective spaces. Show
  that they have negative curvature and are quarter-pinched.
\end{remark}

\begin{remark}[the monodromy map of a locally symmetric space]
  \label{rem:pet-ch10-exercise-lie-algebra-locally-symmetric}
  \lean{PetersenLib.exercise10_5_23}
  \source{petersen:page-0407}
  \uses{rem:pet-ch10-symmetric-space-from-involution}
  \dcref{ch10:10.5.23}
  Give a Lie algebra description of a locally symmetric space (not necessarily
  complete), and explain why it corresponds to a global symmetric space.
  Conclude that a simply connected locally symmetric space admits a monodromy
  map into a unique global symmetric space, bijective when the locally
  symmetric space is complete.
\end{remark}

\begin{remark}[strict curvature-operator sign forces constant curvature]
  \label{rem:pet-ch10-exercise-pinched-curvature-operator-constant}
  \lean{PetersenLib.exercise10_5_24}
  \source{petersen:page-0407}
  \uses{cor:pet-ch10-irreducible-symmetric-einstein}
  \dcref{ch10:10.5.24}
  Show that if an irreducible symmetric space has strictly positive or strictly
  negative curvature operator, then it has constant curvature.
\end{remark}

\begin{remark}[isotropy preserves the translational summand]
  \label{rem:pet-ch10-exercise-bracket-tp-isop}
  \lean{PetersenLib.exercise10_5_25}
  \source{petersen:page-0407}
  \uses{prop:pet-ch10-lie-bracket-curvature}
  \dcref{ch10:10.5.25}
  Let $M$ be a symmetric space. Show that if $X\in t_p$ and $Y\in\mathrm{iso}_p$,
  then $[X,Y]\in t_p$.
\end{remark}

\begin{remark}[curvature and Ricci tensor via the $L_X$ operators]
  \label{rem:pet-ch10-exercise-curvature-ricci-bracket-formula}
  \lean{PetersenLib.exercise10_5_26}
  \source{petersen:page-0407}
  \uses{thm:pet-ch10-curvature-bracket-formula}
  \dcref{ch10:10.5.26}
  Let $M$ be a symmetric space and $X,Y,Z\in t_p$. Show
  $R(X,Y)Z=[L_X,L_Y]Z$ and $\operatorname{Ric}(X,Y)=-\operatorname{tr}([L_X,L_Y])$, where
  $L_X:=\operatorname{ad}_X$.
\end{remark}

\begin{remark}[parallel extension of a $p$-form via holonomy invariance]
  \label{rem:pet-ch10-exercise-parallel-form-holonomy-invariant}
  \lean{PetersenLib.exercise10_5_27}
  \source{petersen:page-0407}
  \uses{rem:pet-ch10-holonomy-properties}
  \dcref{ch10:10.5.27}
  Consider a Riemannian manifold $M$ and a $p$-form $\omega$ on $T_pM$. Show
  that $\omega$ extends to a parallel form on $M$ if and only if $\omega$ is
  invariant under $\mathrm{Hol}_p(M)$.
\end{remark}
